\documentclass[twoside]{article}

\usepackage[preprint]{aistats2027}
\usepackage[round,authoryear]{natbib}
\usepackage{amsmath,amssymb,amsthm}
\usepackage{mathtools}
\usepackage{microtype}
\usepackage{url}

\renewcommand{\bibname}{References}
\renewcommand{\bibsection}{\subsubsection*{\bibname}}
\bibliographystyle{apalike}

\newtheorem{theorem}{Theorem}
\newtheorem{lemma}[theorem]{Lemma}
\newtheorem{proposition}[theorem]{Proposition}
\newcommand{\R}{\mathbb R}
\newcommand{\E}{\mathbb E}
\newcommand{\conv}{\operatorname{conv}}

\begin{document}
\runningauthor{}

\twocolumn[
\aistatstitle{Supplement to ``Sharp Pair Selection for Mean Regression''}

\aistatsauthor{Pahan Dewasurendra}

\aistatsaddress{Johns Hopkins University}
]

\appendix

\section{A Standalone Atomic Proof}\label{app:atomic}

We restate the geometric result in a form that permits repeated support
points and zero masses.

\begin{theorem}\label{thm:app-atomic}
Let $p\in\R^m_+$ satisfy $\sum_i p_i=1$, and let $x_1,\ldots,x_m$ lie in a
real Hilbert space. If $\sum_i p_i x_i=0$, then some $i,j$ satisfy
\[
 \left\|\frac{x_i+x_j}{2}\right\|^2
 \leq
 \max\left\{\frac13,\frac{m-2}{2(m-1)}\right\}
 \sum_i p_i\|x_i\|^2.
\]
Zero-mass points may be removed. The constant is sharp when at most $m$
positive masses are allowed.
\end{theorem}

\paragraph{Proof.}
Remove zero masses and let $m'$ be the remaining support size. The claimed
constant is nondecreasing in support size, so it is enough to prove the result
with $m=m'$. Scale the second moment to one.

If $p_j\geq1/4$, assign probability $(4p_j-1)/3$ to $(j,j)$ and probability
$4p_i/3$ to $(i,j)$ for each $i\ne j$. These probabilities sum to one. Direct
expansion gives
\begin{align*}
 3\E\left\|\frac{x_I+x_J}{2}\right\|^2
 &=(4p_j-1)\|x_j\|^2
   +\sum_{i\ne j}p_i\|x_i+x_j\|^2\\
 &=\sum_i p_i\|x_i\|^2
   +4p_j\|x_j\|^2-\|x_j\|^2\\
 &\quad +(1-4p_j)\|x_j\|^2=1.
\end{align*}
The middle line uses $\sum_{i\ne j}p_i x_i=-p_jx_j$.

Suppose next that every $p_i<1/4$. Put
\begin{align*}
 S&=\sum_i\frac{p_i}{1-4p_i},
 &C&=\frac{S+1}{3S+1},\\
 a_i&=\frac{Cp_i}{(1+S)(1-4p_i)}.
\end{align*}
The identity
\[
 \sum_i\frac{p_i^2}{1-4p_i}=\frac{S-1}{4}
\]
implies
\[
 \sum_i a_i=\frac{CS}{1+S},
 \qquad
 \sum_i p_i a_i=\frac{C(S-1)}{4(1+S)}.
\]
For $i<j$, let $w_{ij}=4(p_i a_j+p_j a_i)$. Then
\[
 \sum_{i<j}w_{ij}
 =4\left(\sum_i a_i-\sum_i p_i a_i\right)
 =\frac{C(3S+1)}{S+1}=1.
\]

Let $G_{ij}=\langle x_i,x_j\rangle$. If
$q=(e_i+e_j)/2$ with probability $w_{ij}$ and $Q=\E[qq^\top]$, then
\[
 Q_{ij}=p_i a_j+a_i p_j\quad(i\ne j).
\]
Also,
\begin{align*}
 Q_{ii}
 &=p_i\sum_{j\ne i}a_j+a_i\sum_{j\ne i}p_j\\
 &=p_iA+a_i(1-2p_i)
 =Cp_i+2p_i a_i,
\end{align*}
where $A=\sum_i a_i$ and the last equality follows from
$a_i(1-4p_i)=p_i(C-A)$. Thus
\[
 Q=C\operatorname{diag}(p)+pa^\top+ap^\top.
\]
Since $Gp=0$,
\[
 \E\left\|\frac{x_I+x_J}{2}\right\|^2
 =\operatorname{tr}(GQ)
 =C\sum_i p_i\|x_i\|^2=C.
\]
Finally, convexity of $t/(1-4t)$ gives
\[
 S\geq m\frac{1/m}{1-4/m}=\frac{m}{m-4}.
\]
Because $(S+1)/(3S+1)$ decreases in $S$, this implies
$C\leq(m-2)/(2(m-1))$. At least one pair is no worse than its expected
cost.

For sharpness, the two-point distribution with masses $3/4$ and $1/4$ gives
$1/3$. A uniform regular $(m-1)$-simplex gives
$(m-2)/(2(m-1))$. Taking the larger proves the result. \hfill$\square$

\section{Matrix Certificate}\label{app:matrix}

The pair distributions can be checked without reference to a particular
geometry. Let $p$ be a positive probability vector and
$\Sigma_p=\operatorname{diag}(p)-pp^\top$. An empirical pair frequency is a
vector $q$ of the form $e_i$ or $(e_i+e_j)/2$.

\begin{proposition}\label{prop:matrix}
In the dominant case, the star distribution satisfies
\[
 \E[(q-p)(q-p)^\top]=\frac13\Sigma_p.
\]
In the diffuse case, the complete-graph distribution satisfies, on the
quotient by the centered direction, the exact constant
$C=(S+1)/(3S+1)$.
\end{proposition}

\paragraph{Proof.}
For the star distribution centered at $j$, direct calculation gives
\[
 \E q=\frac23p+\frac13e_j.
\]
Expanding the second moment then yields the first identity. For the diffuse
distribution, the main paper proves
\[
 \E[qq^\top]
 =C\operatorname{diag}(p)+pa^\top+ap^\top.
\]
If $G$ is any positive semidefinite Gram matrix with $Gp=0$, the last two
terms vanish under the trace pairing. Also
$\operatorname{tr}(Gpp^\top)=0$. Hence
\[
 \operatorname{tr}\left(G\E[(q-p)(q-p)^\top]\right)
 =C\operatorname{tr}(G\Sigma_p).
\]
This is the quotient identity. \hfill$\square$

The star formula is a Loewner-order certificate. The diffuse formula is
tailored to centered geometries, which is exactly the class required after
the mean is translated to zero.

\section{Reduction to Intrinsic Dimension}\label{app:reduction}

For completeness, let $D=\{z_1,\ldots,z_N\}\subset\R^d$ have mean $\mu$.
Translate to $u_i=z_i-\mu$ and let their span have dimension $r$. Consider
the linear program
\begin{align*}
 \min_{p\geq0}\quad&\sum_i p_i\|u_i\|^2\\
 \text{subject to}\quad&\sum_i p_i=1,
 \qquad \sum_i p_i u_i=0.
\end{align*}
The uniform vector is feasible. A basic optimal solution has at most $r+1$
positive entries. Its objective is no larger than the full empirical
variance. Applying Theorem~\ref{thm:app-atomic} to those entries produces a
pair from $D$ with centered squared midpoint at most
\[
 \max\left\{\frac13,\frac{r-1}{2r}\right\}
 \frac1N\sum_i\|z_i-\mu\|^2.
\]
The bias-variance identity converts this into the risk ratio in the main
paper.

\section{Equality for the Diffuse Extremizer}\label{app:equality}

Assume $m\geq5$ and a distribution on exactly $m$ atoms attains
$C_m=(m-2)/(2(m-1))$. It cannot have a mass at least $1/4$, since the star
certificate would give the strictly smaller constant $1/3$. Jensen's
inequality in the diffuse proof must be tight, so $p_i=1/m$ for every $i$.
The complete-graph pair distribution is then uniform over all distinct pairs.

Normalize the variance to one. Every distinct pair must attain the common
minimum $C_m$, since their average equals $C_m$. Let
$b_i=\|x_i\|^2$. For $i\ne j$,
\[
 \langle x_i,x_j\rangle=2C_m-\frac{b_i+b_j}{2}.
\]
Using $\sum_jx_j=0$ and $\sum_jb_j=m$ gives
\[
 0=\left\langle x_i,\sum_jx_j\right\rangle
 =\frac{m-4}{2}(1-b_i).
\]
Thus $b_i=1$ and every off-diagonal inner product is $-1/(m-1)$. The Gram
matrix is that of a regular simplex. This proves the equality statement in
the main paper.

\section{Qualitative Stability}\label{app:stability}

We justify the stability statement in the main paper. Fix $m\geq5$ and
consider a sequence of diffuse mass vectors $p^{(k)}$ and centered,
unit-variance Gram matrices $G^{(k)}$. Suppose their best pair costs converge
to $C_m=(m-2)/(2(m-1))$.

The exact diffuse certificate gives
\[
 \min_{i,j}\frac{G^{(k)}_{ii}+G^{(k)}_{jj}+2G^{(k)}_{ij}}4
 \leq C(p^{(k)})\leq C_m.
\]
The left side converges to $C_m$, so $C(p^{(k)})\to C_m$. Strict convexity in
the Jensen step forces $p^{(k)}\to(1/m,\ldots,1/m)$, after a fixed labeling.
The masses are therefore bounded away from zero. Unit weighted trace then
bounds every diagonal entry of $G^{(k)}$. Positive semidefiniteness and
Cauchy--Schwarz bound every off-diagonal entry, so the Gram matrices are
precompact.

For the diffuse certificate, every distinct-pair weight converges to
$2/(m(m-1))$. Its weighted average pair cost equals $C(p^{(k)})$, while every
pair cost is bounded below by the best pair cost. Hence every distinct-pair
cost converges to $C_m$. Any subsequential Gram limit is centered under
uniform masses, has unit weighted trace, and has all distinct midpoint costs
equal to $C_m$. The equality argument in Section~\ref{app:equality} identifies
the limit as the regular-simplex Gram matrix. Thus the full sequence converges
to that matrix, up to relabeling and isometry.

\end{document}
